Empathetic language is not sufficient for mental-health support dialogue. A response can sound caring while skipping therapeutic steps, mixing incompatible techniques, overusing peer advice, or leaving unsafe decisions unauditable. We present \trace, a stage-aware Vietnamese student support prototype built around auditable process control rather than free-form response generation. \trace\ maintains a dynamic blackboard state over user concerns, distress, thoughts, route evidence, stage milestones, safety flags, delivered interventions, and peer history. This state controls route pacing, optional Yalom-style peer contributions, validator repair, and post-generation safety review. We also define VSM, a session-level benchmark protocol for measuring route fidelity, stage fidelity, technique fidelity, safety, peer dynamics, reliability, and latency across \trace, ablations, and prompt baselines. The paper provides the system design and frozen evaluation plan; quantitative claims are reserved for the full VSM core run, LLM judge pass, and human audit.
